\documentclass[10pt,a4paper]{amsart}
\usepackage[margin=19mm]{geometry}
\usepackage{amsmath,amssymb,mathtools}
\usepackage[scr=rsfs,cal=euler]{mathalfa}
\usepackage{xfrac}
\usepackage[unicode,hidelinks,bookmarks=false]{hyperref}
\newcommand{\ie}{i.e.\ }
\newcommand{\cf}{cf.\ }
\newcommand{\resp}{respectively\ }
\newcommand{\Subsection}{\subsection}
\providecommand{\nobreakdash}{\nobreak\hbox{-}}
\setlength{\emergencystretch}{2em}
\multlinegap0pt
\usepackage{amssymb}
\usepackage{float}
\newtheorem{theorem}{Theorem}
\title{Moduli difference of inverse logarithmic coefficients of univalent functions}
\author{Vasudevarao Allu, Amal Shaji}
\date{Source: arXiv:2403.10031v1 (CC BY 4.0)}
\begin{document}
\maketitle

\noindent\emph{Source and licence.} Moduli difference of inverse logarithmic coefficients of univalent functions. Version arXiv:2403.10031v1 (CC BY 4.0), distributed by arXiv under the Creative Commons Attribution 4.0 International licence, http://creativecommons.org/licenses/by/4.0/, as recorded in the arXiv licence metadata for this exact version and shown on the abstract page https://arxiv.org/abs/2403.10031v1. Full licence notice: \texttt{licenses/arxiv-2403.10031-CC-BY-4.0.txt}.\par\smallskip

\noindent\textit{Editorial scope.} Counted unit: the article's theorem classS (source0.tex lines 384--399). The article's complete original proof of it is delivered with it (source0.tex lines 400--446, one \texttt{proof} environment of 47 lines). No mathematics is written, repaired or re-derived: the statement and the proof are the article's own lines, in the article's own notation, and no retained line is substituted or re-flowed.  Retained uncounted as the source-local context the proof needs: lines 132--134; lines 228--233.  Nothing was omitted from any retained span, so the package carries no omission record.  No retained line contains a citation, so the wrapper carries no bibliography and no bibliography entry.  No retained line contains a figure environment, an image-inclusion command or drawing code, so no image is delivered and the package declares no asset dependency.  Editorial formatting only, in addition: an AMS \texttt{amsart} preamble that restates the article's own declarations and macros used by the retained lines, namely the environments \texttt{theorem} on separate counters, together with the standard packages those lines need.  The retained environments print in this document as Theorem 1 in order of appearance, whereas the article prints the same items with the numbering of its own full document.  The complete native source of the article as distributed in the arXiv e-print for this exact version is bundled unchanged as \texttt{sources/156142/source0.tex}.
\begin{equation}\label{S}
	f(z)= z+\sum_{n=2}^{\infty}a_n z^n, \quad z\in \mathbb{D}.
\end{equation}

\begin{equation}\label{Gammaaa}
\begin{aligned}
& \Gamma_{1}=-\cfrac{1}{2}\,a_2,\\[2mm]
& \Gamma_2=-\cfrac{1}{2}\, a_3+\cfrac{3}{4}\, a_2^2.
\end{aligned}
\end{equation}

 \begin{theorem}\label{classS}
 Let $f \in \mathcal{S}$ of the form \eqref{S}, then 
 $$
 |\Gamma_2|-|\Gamma_1| \leq \cfrac{1}{2}
 $$
 and 
 \begin{equation*}\label{BBB}
|\Gamma_2|-|\Gamma_1| \geq \left \{
	\begin{array}{ll}
		-\cfrac{1}{2} & {\mbox{ if }} |a_2|\leq 1,
		\\[3mm]
		-0.6353..., & {\mbox{ if }} |a_2|>1.
	\end{array}
	\right.
\end{equation*}
 \end{theorem}

 \begin{proof}
 Let $f \in \mathcal{S}$ of the form \eqref{S}. Then by \eqref{Gammaaa}, we have
\begin{equation}\label{ga}
|\Gamma_2|-|\Gamma_1|=\cfrac{1}{2}\left(|a_3-\cfrac{3}{2}a_2^2|-|a_2|\right).
\end{equation}
From the Bieberbach conjecture, for all $f \in \mathcal{S}$, we have $\cfrac{1}{2}|a_2| \leq 1$, and so by \eqref{ga}

\begin{equation}
\begin{aligned}
2(|\Gamma_2|-|\Gamma_1|)&\leq \left|a_3-\cfrac{3}{2}a_2^2\right|-\cfrac{1}{2}\left|a_2^2\right|\\[2mm]
&\leq \left| a_3-\cfrac{3}{2}a_2^2+\cfrac{1}{2}a_2^2 \right| \\[2mm]
&\leq \left| a_3-a_2^2\right|\\[2mm]
&\leq 1,
\end{aligned}
\end{equation}
and hence we have
$$
|\Gamma_2|-|\Gamma_1| \leq \cfrac{1}{2}
$$

\noindent Now we obtain the lower bound for $|\Gamma_2|-|\Gamma_1|$.
We have 
$$
|\Gamma_1|-|\Gamma_2| = \left(\cfrac{1}{2}\left|a_2\right|-\left|a_3-\cfrac{3}{2}a_2^2\right|\right).
$$
If $|a_2|\leq 1$, then clearly $|\Gamma_1|-|\Gamma_2| \leq 1/2$. If $|a_2|> 1$, then
\begin{equation}\label{fek}
\begin{aligned}
2(|\Gamma_1|-|\Gamma_2|) &= \left|a_2\right|-\left|a_3-\cfrac{3}{2}a_2^2\right|\\
&\leq \left|a_2\right|^2-\left|a_3-\cfrac{3}{2}a_2^2\right|\\
&\leq \left| a_3-\cfrac{3}{2}a_2^2+a_2^2\right|\\
&=\left| a_3-\frac{1}{2}a_2^2\right|
\end{aligned}
\end{equation}
By taking $\mu=1/2$ in Fekete-Szegö Theorem and using \eqref{fek}, we obtain 
$$
|\Gamma_1|-|\Gamma_2|\leq  \cfrac{1+2e^{-2}}{2}=0.6353...
$$
For the upper bound, equality holds for the Koebe function $k(z)=z/(1-z)^2$. For the lower bound, when $|a_2|\leq 1$, equality holds for the function
$$
f(z)=\cfrac{z}{1+z+z^2}.
$$
For $|a_2| > 1$, whether the lower bound $-0.6353...$ is
the best possible is an open problem.


 \end{proof}

\end{document}
